% ---------------------------------------------------------------- page de garde
\begin{titlepage}
\begin{center}
  {\large \textsc{République Tunisienne}}\\[2pt]
  {\small Ministère de l'Enseignement Supérieur et de la Recherche Scientifique}\\[14pt]
  % Logos : déposer les fichiers dans figures/ puis décommenter
  % \includegraphics[height=2cm]{figures/logo_esprit.png}\hfill
  % \includegraphics[height=2cm]{figures/logo_tukhnanutha.png}\\[10pt]
  {\Large \textbf{ESPRIT} --- École Supérieure Privée d'Ingénierie et de Technologies}\\[30pt]

  {\LARGE \textsc{Rapport de projet de fin d'études}}\\[8pt]
  {\large Présenté en vue de l'obtention du\\
  \textbf{Diplôme National d'Ingénieur} \acompleter{spécialité / option}}\\[26pt]

  \rule{\linewidth}{0.6pt}\\[10pt]
  {\Large \bfseries \color{bleuespirt} \titrerapport}\\[8pt]
  \rule{\linewidth}{0.6pt}\\[28pt]

  {\large Réalisé par : \textbf{\autrice}}\\[24pt]

  \begin{tabular}{ll}
    Encadrant académique : & \acompleter{Nom et prénom}\\[4pt]
    Encadrant professionnel : & \acompleter{Nom et prénom (Tukhnanutha)}\\
  \end{tabular}\\[30pt]

  {\large Organisme d'accueil : \textbf{Tukhnanutha}, Bizerte}\\[4pt]
  {Projet ANEO --- application ATOM Wind}\\[40pt]

  \vfill
  {\large Année universitaire 2025--2026}
\end{center}
\end{titlepage}

% ---------------------------------------------------------------- dédicace
\chapter*{Dédicace}
\addcontentsline{toc}{chapter}{Dédicace}
\begin{flushright}
\itshape
\acompleter{Texte de la dédicace, rédigé par l'autrice.}
\end{flushright}

% ---------------------------------------------------------------- remerciements
\chapter*{Remerciements}
\addcontentsline{toc}{chapter}{Remerciements}
\acompleter{Remerciements personnels : encadrant académique, encadrant professionnel et équipe de
Tukhnanutha, membres du jury, famille.}

% ---------------------------------------------------------------- résumé
\chapter*{Résumé}
\addcontentsline{toc}{chapter}{Résumé}
Ce projet de fin d'études, réalisé chez Tukhnanutha dans le cadre de l'écosystème ANEO, porte sur
la reconstruction du champ de vent tridimensionnel (composantes $u$, $v$, $w$) entre 25 et 1\,000\,m
d'altitude à partir de mesures de surface éparses. Nous avons conçu \emph{PC-FNO3D}, un opérateur
neuronal de Fourier 3D contraint par la physique : la perte combine l'erreur sur les données, la
régularité des gradients, la conservation anélastique de la masse $\nabla\cdot(\rho\vec{v})=0$ et
la condition au sol. Sur un jeu synthétique de 100 jours, le modèle de référence (v7) atteint
cinq des six objectifs mesurés (résidu de masse $2{,}6\cdot10^{-5}$, inférence en 91\,ms sur GPU,
robustesse à 25\,\% des capteurs) et divise par 7,5 à 9,8 l'erreur maximale des méthodes de
référence (profil logarithmique ou loi de puissance avec interpolation IDW ou krigeage). Le modèle a
ensuite été transposé au domaine réel du Grand Tunis (68\,$\times$\,56\,km) en apprenant une
correction du profil ERA5, avec une vérité générée forcée par ERA5 et une validation sur les
anémomètres réels et les radiosondages. Le modèle livré (v4) obtient une erreur globale de
0,38\,m/s, un résidu de masse de $1{,}05\cdot10^{-5}$ et une couverture calibrée de 88\,\% pour
l'intervalle à 80\,\%. Le tout est servi par une API sécurisée, un serveur MCP et l'application
web trilingue ATOM Wind (vue 3D, conseils agricoles de J0 à J+2, démonstration éolienne). Les
limites sont rapportées sans les masquer : le niveau le plus bas reste à 0,53\,m/s, et l'étude de
dimensionnement du réseau (OS6) montre que le modèle actuel ne tire aucun gain de capteurs ajoutés.

\noindent\textbf{Mots-clés :} champ de vent 3D, opérateur neuronal de Fourier, apprentissage
contraint par la physique, ERA5, GFS, quantification de l'incertitude, application web.

% ---------------------------------------------------------------- abstract
\chapter*{Abstract}
\addcontentsline{toc}{chapter}{Abstract}
\begin{otherlanguage}{english}
This final-year project, carried out at Tukhnanutha within the ANEO ecosystem, addresses the
reconstruction of the three-dimensional wind field ($u$, $v$, $w$) between 25 and 1,000\,m above
sea level from sparse surface measurements. We designed \emph{PC-FNO3D}, a physics-constrained 3D
Fourier neural operator whose loss combines the data error, a gradient term, anelastic mass
conservation $\nabla\cdot(\rho\vec{v})=0$ and the ground boundary condition. On a 100-day synthetic
dataset the reference model (v7) meets five of the six measured objectives (mass residual
$2.6\cdot10^{-5}$, 91\,ms GPU inference, robustness with 25\,\% of the sensors) and reduces the
worst-level error of the reference methods (log or power-law profile with IDW or kriging) by a
factor of 7.5 to 9.8. The model was then transferred to the real Greater Tunis domain
(68\,$\times$\,56\,km) by learning a correction of the ERA5 profile, with an ERA5-forced generated
truth and validation against real anemometers and radiosondes. The delivered model (v4) reaches a
global error of 0.38\,m/s, a mass residual of $1.05\cdot10^{-5}$ and a calibrated 88\,\% coverage
of the 80\,\% interval. It is served through a secured API, an MCP server and the trilingual ATOM
Wind web application (3D view, farming advice from day 0 to day 2, wind-energy demonstration).
Limits are reported as measured: the lowest level remains at 0.53\,m/s, and the network-sizing
study (OS6) shows that the current model gains nothing from additional sensors.

\noindent\textbf{Keywords:} 3D wind field, Fourier neural operator, physics-constrained learning,
ERA5, GFS, uncertainty quantification, web application.
\end{otherlanguage}
